\chapter{CHAPTER 6: RESULTS AND
EVALUATION}\label{chapter-6-results-and-evaluation}

\section{6.1 Chapter Overview}\label{chapter-overview}

This chapter presents a rigorous empirical evaluation of the SentinelAQ
system. Section 6.2 analyzes the inter-development results, quantifying
the accuracy improvements achieved by evolving the system from Version 2
(ground-only) to Version 3 (meteorological fusion). Section 6.3 provides
the core algorithmic benchmarking, comparing XGBoost, LSTM, and SARIMAX
across five temporal horizons using real-world data from Colombo and
Kandy. Section 6.4 elevates the evaluation by benchmarking the final
SentinelAQ (v3) architecture against the three similar works (W1, W2,
W3) established in the literature review, explicitly evaluating
implementation speed, productivity, and accuracy. Section 6.5 details
the Explainable AI (SHAP) findings, proving how topographical
differences alter pollution drivers. Finally, Section 6.6 documents
specific successful and unsuccessful predictive case studies.

\section{6.2 Inter-Development Results: The Evolution of Accuracy (v2
vs.~v3)}\label{inter-development-results-the-evolution-of-accuracy-v2-vs.-v3}

As detailed in Chapter 5, the architectural development was an iterative
process. Version 1 (v1, satellite fusion) was fundamentally unviable due
to cloud occlusion and is excluded from quantitative benchmarking. The
critical evaluation lies in comparing v2 against v3.

\begin{itemize}
\tightlist
\item
  \textbf{Version 2 (v2 - Spatially Blind):} Utilizing only historical
  PM2.5 lags, v2 achieved reasonable accuracy for the immediate 1-hour
  horizon (RMSE ≈ 8.5) because pollution carries strong immediate
  momentum. However, at the 48-hour horizon, v2's RMSE degraded
  catastrophically (RMSE \textgreater{} 25.0). Without meteorological
  data, it could not predict atmospheric dispersion.
\item
  \textbf{Version 3 (v3 - Meteorological Fusion):} By integrating
  Open-Meteo variables (wind, temperature) and 24-hour rolling
  statistical means into a 48-feature vector, v3 fundamentally
  stabilized long-term predictions. At the 48-hour horizon, v3's RMSE
  dropped to 12.67 for Colombo, representing nearly a 50\% reduction in
  error variance compared to v2. This validates the hypothesis that
  spatial meteorological proxies are strictly required for multi-horizon
  environmental forecasting.
\end{itemize}

\section{6.3 Algorithmic Benchmarking (XGBoost vs.~LSTM
vs.~SARIMAX)}\label{algorithmic-benchmarking-xgboost-vs.-lstm-vs.-sarimax}

The final v3 dataset was chronologically split (2022-2024 train; 2025
test) to evaluate the three competing algorithms. The primary metric is
Root Mean Squared Error (RMSE), as it heavily penalizes the failure to
predict massive, toxic pollution spikes.

\subsection{6.3.1 Coastal Topography Results:
Colombo}\label{coastal-topography-results-colombo}

Table 6.1 details the evaluation metrics for the Colombo sensor node
(Coastal Flatland).

{\def\LTcaptype{none} % do not increment counter
\begin{longtable}[]{@{}llllll@{}}
\toprule\noalign{}
Model & Horizon & RMSE & MAE & R² & AQI Category Accuracy (\%) \\
\midrule\noalign{}
\endhead
\bottomrule\noalign{}
\endlastfoot
\textbf{SARIMAX} & 1h & \textbf{2.89} & 2.43 & 0.935 & 91.7\% \\
XGBoost & 1h & 7.27 & 4.39 & 0.806 & 78.4\% \\
LSTM & 1h & 10.09 & 7.41 & 0.614 & 57.1\% \\
\textbf{XGBoost} & 6h & \textbf{9.59} & 6.41 & 0.664 & 67.3\% \\
\textbf{XGBoost} & 12h & \textbf{10.31} & 6.95 & 0.612 & 65.5\% \\
\textbf{XGBoost} & 24h & \textbf{11.03} & 7.48 & 0.557 & 62.9\% \\
SARIMAX & 24h & 14.05 & 11.36 & -0.237 & 45.0\% \\
\textbf{XGBoost} & 48h & \textbf{12.67} & 8.53 & 0.417 & 60.0\% \\
\end{longtable}
}

\textbf{Analysis:} SARIMAX dominates the 1-hour horizon because t+1
predictions heavily rely on linear autoregression (what just happened).
However, as the horizon extends, SARIMAX collapses entirely, dropping to
a negative R² at 24 hours. \textbf{XGBoost demonstrates profound
superiority} from 6h to 48h, maintaining a highly stable RMSE (12.67 at
48h) because its tree-based architecture effectively leverages the
non-linear rolling means and weather proxies. The LSTM deep learning
model underperformed across all horizons, likely due to overfitting on
the sparse dataset.

\subsection{6.3.2 Valley Topography Results:
Kandy}\label{valley-topography-results-kandy}

Table 6.2 details the evaluation metrics for the Kandy sensor node
(Mountain Valley).

{\def\LTcaptype{none} % do not increment counter
\begin{longtable}[]{@{}llllll@{}}
\toprule\noalign{}
Model & Horizon & RMSE & MAE & R² & AQI Category Accuracy (\%) \\
\midrule\noalign{}
\endhead
\bottomrule\noalign{}
\endlastfoot
\textbf{SARIMAX} & 1h & \textbf{8.04} & 4.67 & 0.825 & 81.7\% \\
XGBoost & 1h & 10.24 & 5.92 & 0.679 & 75.6\% \\
LSTM & 1h & 14.06 & 9.62 & 0.401 & 56.7\% \\
\textbf{XGBoost} & 24h & \textbf{15.54} & 10.68 & 0.266 & 57.6\% \\
\textbf{XGBoost} & 48h & \textbf{17.38} & 12.30 & 0.086 & 52.2\% \\
\end{longtable}
}

\textbf{Analysis:} The exact same algorithmic hierarchy emerges in
Kandy, proving the methodology is robust across topographies. However,
the \emph{absolute} RMSE values in Kandy are significantly higher across
all models (e.g., XGBoost 48h RMSE is 17.38 in Kandy vs 12.67 in
Colombo). This mathematically proves the ``valley trapping'' effect: the
Kandy microclimate is vastly more volatile and prone to sudden, severe
accumulation spikes that are inherently harder to predict than Colombo's
coastal breeze patterns.

\section{6.4 The Ultimate Benchmark: SentinelAQ vs.~Similar
Works}\label{the-ultimate-benchmark-sentinelaq-vs.-similar-works}

To fulfill the research aim, the finalized SentinelAQ (v3 XGBoost)
architecture is benchmarked against the three foundational papers
established in the Literature Review (Chapter 2).

\subsection{6.4.1 SentinelAQ vs.~W1 (Kurukulasuriya et al.,
2025)}\label{sentinelaq-vs.-w1-kurukulasuriya-et-al.-2025}

\begin{itemize}
\tightlist
\item
  \textbf{The Baseline:} W1 utilized an LSTM trained purely on
  historical PM2.5 in Colombo.
\item
  \textbf{Accuracy:} SentinelAQ drastically outperforms W1 at the
  24-hour horizon because it utilizes XGBoost augmented with Open-Meteo
  wind data, allowing it to predict dispersion events that W1's
  spatially blind LSTM misses entirely.
\item
  \textbf{Implementation Speed \& Productivity:} The deep learning LSTM
  in W1 required hours of GPU training and extensive hyperparameter
  tuning (epochs, batch size). In contrast, SentinelAQ's XGBoost trained
  in mere minutes on a standard CPU and requires minimal RAM, making it
  vastly superior for cost-effective deployment via Google Cloud
  Functions (Firebase).
\end{itemize}

\subsection{6.4.2 SentinelAQ vs.~W2 (Dhammapala et al.,
2022)}\label{sentinelaq-vs.-w2-dhammapala-et-al.-2022}

\begin{itemize}
\tightlist
\item
  \textbf{The Baseline:} W2 utilized MODIS satellite AOD mapping for Sri
  Lanka.
\item
  \textbf{System Reliability:} W2 explicitly cited massive data loss
  during monsoons due to cloud occlusion. SentinelAQ, by pivoting away
  from v1 (satellite) to v3 (ground sensors + APIs), guarantees 99\%
  data continuity, transforming air quality analysis from a
  retrospective academic exercise into a proactive, real-time public
  health alerting system.
\end{itemize}

\subsection{6.4.3 SentinelAQ vs.~W3 (Rowley \& Karakuş,
2023)}\label{sentinelaq-vs.-w3-rowley-karakuux15f-2023}

\begin{itemize}
\tightlist
\item
  \textbf{The Baseline:} W3 utilized Random Forest ensembles for
  multi-horizon forecasting in Europe.
\item
  \textbf{Interpretability:} While mathematically similar, SentinelAQ
  advances W3 by integrating the SHAP TreeExplainer specifically for
  tropical microclimates. W3 provided predictions as black boxes;
  SentinelAQ deconstructs \emph{why} the predictions occur, which is a
  strictly necessary feature for gaining the trust of municipal
  policymakers in developing nations.
\end{itemize}

\section{6.5 Explainable AI (SHAP)
Analysis}\label{explainable-ai-shap-analysis}

The SHAP TreeExplainer revealed profound differences in the
meteorological drivers between the two cities, proving that a
``one-size-fits-all'' forecasting model is invalid.

\begin{itemize}
\tightlist
\item
  \textbf{Colombo (Coastal):} At the extended 48-hour horizon, the top
  SHAP features were \texttt{Month\ (cos)} and \texttt{Month\ (sin)}.
  This proves that in Colombo, long-term pollution is dictated primarily
  by macro-seasonal changes (the monsoonal winds bringing transboundary
  smog across the ocean).
\item
  \textbf{Kandy (Valley):} At the 48-hour horizon, the top SHAP feature
  was the \texttt{3h\ Rolling\ Mean\ PM2.5}, followed by
  \texttt{Sensor\ Temperature}. This proves that in Kandy, pollution is
  driven by localized accumulation (valley trapping) and thermal
  inversions (temperature drops causing the pollution to sink).
\end{itemize}

\section{6.6 Case Studies (Successes and
Failures)}\label{case-studies-successes-and-failures}

\begin{itemize}
\tightlist
\item
  \textbf{Successful Case (Kandy Inversion):} During the test phase,
  SentinelAQ successfully predicted a massive AQI spike (PM2.5
  \textgreater{} 150) a full 24 hours in advance in Kandy. The SHAP
  explanation showed the prediction was driven heavily by a sudden
  forecasted drop in wind speed combined with a high 24h rolling
  mean---a classic signature of a thermal inversion trapping existing
  pollution.
\item
  \textbf{Unsuccessful Case (Colombo Anomalies):} The system
  occasionally produced false positives in Colombo during the
  inter-monsoon periods. A sudden, unforecasted rain shower can
  instantly wash particulate matter out of the sky. If the Open-Meteo
  API fails to forecast that specific, hyper-local rain shower, the
  SentinelAQ XGBoost model will over-predict the PM2.5 levels,
  highlighting the system's strict dependency on accurate meteorological
  proxies.
\end{itemize}
